% Plik generowany: npm --prefix 04_app run gen:report — nie edytować ręcznie.
\begingroup\small
\begin{longtable}{@{}p{0.16\linewidth} p{0.4\linewidth} p{0.36\linewidth}@{}}
\caption{Najefektywniejsze zastosowanie i główne słabości mechanizmów.}\label{tab:zastosowanie}\\
\toprule
Mechanizm & Gdzie i po co & Główne słabości \\
\midrule
\endhead
Hasło & Jako jeden z dwóch czynników albo w serwisach niskiego ryzyka; z menedżerem haseł i unikalnym hasłem na serwis. & Phishing, powtarzanie haseł w wielu serwisach, słabe hasła łamane po wycieku bazy. \\
SMS OTP & Jako przejściowy drugi składnik tam, gdzie nie ma nic lepszego; nie do kont o wysokiej wartości. & SIM-swap, phishing w czasie rzeczywistym, przechwycenie w sieci (SS7), podgląd na zablokowanym ekranie. \\
TOTP (aplikacja) & Jako drugi składnik zamiast SMS: działa offline i nie zależy od operatora. & Phishing w czasie rzeczywistym, kradzież sekretu (np. z niezaszyfrowanej kopii zapasowej), utrata telefonu bez kopii. \\
Push + number matching & Firmy i banki z własną aplikacją; zawsze z dopasowaniem liczby i kontekstem (kto, skąd, co). & Zmęczenie powiadomieniami (bez dopasowania liczby), socjotechnika przez telefon, phishing w czasie rzeczywistym. \\
Klucz sprzętowy FIDO2 & Konta uprzywilejowane i pracownicy firm (B2B), administratorzy, osoby szczególnie narażone. & Koszt i logistyka (zakup, zapasowy klucz), zgubienie, słabsze ścieżki odzyskiwania konta. \\
Passkey & Logowanie klientów (B2C) zamiast hasła i SMS; w firmie w wersji związanej z urządzeniem. & Odzyskiwanie konta, przejęcie konta chmurowego (passkey synchronizowany), kradzież sesji po zalogowaniu. \\
Biometria & Lokalne odblokowanie passkeya lub aplikacji banku; wygoda bez wysyłania danych biometrycznych na serwer. & Ataki prezentacji (sztuczny palec, zdjęcie), nie da się jej zmienić po wycieku, błędy dopasowania (FAR/FRR). \\
Karta płatnicza & Płatności w sklepie (chip, zbliżeniowo); pasek magnetyczny tylko jako przestarzała rezerwa. & Skimming i klonowanie paska, przekaźnik NFC (relay), obejścia PIN w niektórych konfiguracjach, phishing danych karty online. \\
MFA i ocena ryzyka & Wszędzie, gdzie stawka jest wysoka; z odpornym na phishing składnikiem i bez słabych ścieżek zapasowych. & Słabe metody zapasowe (SMS, pytania), zmęczenie użytkownika, kradzież sesji po zalogowaniu. \\
BLIK & Płatności w internecie i wypłaty w Polsce (C2C, B2C) bez podawania danych karty. & Socjotechnika: „pożycz na BLIKa”, fałszywe „odbieranie płatności”, zatwierdzanie bez czytania. \\
\bottomrule
\end{longtable}
\endgroup
